% Fig 1 — the L0->L3 stack map, placing Paper 3 (the bridge). Harbor Volume idiom.
% Design: one bold title + small-caps whom per rung, one italic content line with
% an honest-state tag; Paper 3 is a cobalt bracket spanning the L2->L3 gap.
% The through-line chain lives in the caption.
\providecolor{hhsand}{HTML}{E9DCC4}
\providecolor{hhsanddeep}{HTML}{D8C7A6}
\providecolor{hhebony}{HTML}{121212}
\providecolor{hhink}{HTML}{1B1712}
\providecolor{hhcobalt}{HTML}{003FB8}
\providecolor{hhamber}{HTML}{6B4500}
\providecolor{hhteal}{HTML}{00564C}
\providecolor{hhpaper}{HTML}{FBF7EF}

\begin{figure}[H]
\centering
\resizebox{0.95\textwidth}{!}{%
\begin{tikzpicture}[
  font=\small,
  rung/.style={draw=hhebony,thick,rounded corners=3pt,inner sep=7pt,
               minimum width=132mm,minimum height=15mm,align=left,text width=124mm},
  paper/.style={draw=hhink,thick,fill=hhpaper,rounded corners=2pt,inner sep=4pt,
                font=\footnotesize,align=center,text width=22mm},
]
% The four rungs, L0 bottom -> L3 top
\node[rung,fill=hhsand!30] (l0) at (0,0)
  {{\bfseries L0 --- Daemon (kernel)}\hfill{\footnotesize\textsc{for the machine}}\\[2pt]
   \textit{always-on local SQLite truth: ports, claims, sessions, tube, memory}\hfill{\footnotesize\textsc{built}}};
\node[rung,fill=hhsand!30,above=5mm of l0] (l1)
  {{\bfseries L1 --- Coordination protocol}\hfill{\footnotesize\textsc{for the agents}}\\[2pt]
   \textit{typed conversation, commitments, delegation, Arbiter}\hfill{\footnotesize\textsc{designed}}};
\node[rung,fill=hhsand!30,above=5mm of l1] (l2)
  {{\bfseries L2 --- Legibility \& authority (the wedge)}\hfill{\footnotesize\textsc{for the operator}}\\[2pt]
   \textit{digest-with-zoom, roadmap-as-truth, adversarial review, the console}\hfill{\footnotesize\textsc{built}}};
\node[rung,fill=hhsand!30,above=5mm of l2] (l3)
  {{\bfseries L3 --- Economy \& federation}\hfill{\footnotesize\textsc{for the market}}\\[2pt]
   \textit{anchor, escrow, reputation, federation, the labor/skill market}\hfill{\footnotesize\textsc{whitepaper'd}}};

% Paper 3 = the bridge from L2 to L3: a cobalt bracket spanning the gap
\draw[hhcobalt,very thick] ([xshift=-4mm]l2.north west) -- ([xshift=-7mm]l2.north west)
  -- ([xshift=-7mm]l3.south west) -- ([xshift=-4mm]l3.south west);
\node[paper,draw=hhcobalt,very thick,anchor=east,align=center]
  at ([xshift=-9mm]$(l2.north west)!0.5!(l3.south west)$)
  {\textbf{Paper 3}\\\emph{From Spawn to Person}\\\textit{the bridge}};

% companion papers, left column
\node[paper,anchor=east] at ([xshift=-9mm]l0.west) {Paper 2\\\emph{Anchor} (L0--L1)};
\node[paper,anchor=east] at ([xshift=-9mm]l2.west) {Paper 1\\\emph{Legible Swarm}};
\node[paper,anchor=east] at ([xshift=-9mm]l3.west) {Paper 4\\\emph{Harbor Economy}};
\end{tikzpicture}}
\caption{\textbf{The Harbor Volume stack, and where this paper sits.} Each rung is
for a different \emph{whom}; the honest state of each is labelled. This paper is the
\emph{bridge} from L2 (a legible \emph{person} doing work) to L3 (a tradeable
\emph{reputation}) --- the through-line memory $\to$ checkpoint $\to$
\emph{continuity} $\to$ a \emph{person} $\to$ outcomes $\to$ \emph{reputation} $\to$
the market. Paper~2 supplies the cryptographic identity below; Paper~4 consumes the
reputation substrate above.}
\label{fig:stack}
\end{figure}
